\par\Needspace{6\baselineskip}\medskip
\hypertarget{native-146-article106}{}\pdfbookmark[2]{106.}{article-106}
\begin{center}106.\end{center}
In tractatu de relationibus maxime insignibus ad motum corporum coelestium in sectionibus conicis spectantibus silentio praeterire non possumus expressionem elegantem temporis per semiaxem maiorem, summam $r+r'$ atque chordam duo loca iungentem. Haec formula pro parabola quidem primo \SourceLeaf{147}ab ill. \textsc{Euler} inventa esse videtur (Miscell. Berolin. T. VII p. 20), qui tamen eam in posterum neglexit, neque etiam ad ellipsin et hyperbolam extendit: errant itaque, qui formulam clar. \textsc{Lambert} tribuunt, etiamsi huic geometrae meritum, hanc expressionem oblivione sepultam proprio marte eruisse et ad reliquas sectiones conicas ampliavisse, non possit denegari. Quamquam hoc argumentum a pluribus geometris iam tractatum sit, tamen lectores attenti expositionem sequentem haud superfluam agnoscent. A motu elliptico initium facimus.
\par
Ante omnia observamus, angulum circa Solem descriptum $2f$ (art. 88, unde reliqua quoque signa desumimus) infra $360^0$ supponi posse; patet enim, si iste angulus 360 gradibus augeatur, tempus una revolutione sive $\dfrac{a^{\frac32}\mathpunct{.}360^0}{k}=a^{\frac32}\times365{,}25$ diebus crescere. Iam si chordam per $\rho$ denotamus, manifestum est fieri
\[
\rho\rho=(r'\cos v'-r\cos v)^2+(r'\sin v'-r\sin v)^2
\]
\noindent adeoque per aequationes VIII, IX art. 8
\[
\begin{aligned}
\rho\rho&=aa(\cos E'-\cos E)^2+aa\cos\varphi^2(\sin E'-\sin E)^2\\[.6em]
&=4aa\sin g^2(\sin G^2+\cos\varphi^2\cos G^2)\\[.6em]
&=4aa\sin g^2(1-ee\cos G^2).
\end{aligned}
\]
\noindent Introducamus angulum auxiliarem $h$ talem, ut sit $\cos h=e\cos G$; simul, quo omnis ambiguitas tollatur, supponemus, $h$ accipi inter 0 et $180^0$, unde $\sin h$ erit quantitas positiva. Quoniam itaque etiam $g$ inter eosdem limites iacet (si enim $2g$ ad $360^0$ vel ultra ascenderet, motus circa Solem revolutionem integram attingeret vel superaret), ex aequatione praecedente sponte sequitur
\[
\rho=2a\sin g\sin h,
\]
\noindent siquidem chorda tamquam quantitas positiva consideratur. Quum porro habeatur
\[
r+r'=2a(1-e\cos g\cos G)=2a(1-\cos g\cos h),
\]
\noindent patet, si statuatur $h-g=\delta$, $h+g=\epsilon$, fieri
\[
\begin{aligned}
\mathrm{1)}\qquad&r+r'-\rho=2a(1-\cos\delta)=4a\sin\tfrac12\delta^2\\[.6em]
\mathrm{2)}\qquad&r+r'+\rho=2a(1-\cos\epsilon)=4a\sin\tfrac12\epsilon^2.
\end{aligned}
\]
\par
Denique habetur
\[
kt=a^{\frac32}(2g-2e\sin g\cos G)=a^{\frac32}(2g-2\sin g\cos h),
\]
\noindent\SourceLeaf{148}sive
\[
\mathrm{3)}\qquad kt=a^{\frac32}(\epsilon-\sin\epsilon-(\delta-\sin\delta)).
\]
\par
Determinari poterunt itaque, secundum aequationes 1, 2, anguli $\delta$ et $\epsilon$ ex $r+r'$, $\rho$ et $a$: quamobrem ex iisdem quantitatibus determinabitur, secundum aequationem 3, tempus $t$. Si magis placet, haec formula ita exhiberi potest:
\[
\begin{aligned}
kt=a^{\frac32}\Biggl\{&\operatorname{arc}\cos\frac{2a-(r+r')-\rho}{2a}\\[.6em]
&-\sin\operatorname{arc}\cos\frac{2a-(r+r')-\rho}{2a}\\[.6em]
&-\operatorname{arc}\cos\frac{2a-(r+r')+\rho}{2a}\\[.6em]
&+\sin\operatorname{arc}\cos\frac{2a-(r+r')+\rho}{2a}\Biggr\}\mathpunct{\cdot}
\end{aligned}
\]
\noindent Sed in determinatione angulorum $\delta$, $\epsilon$ per cosinus suos ambiguitas remanet, quam propius considerare oportet. Sponte quidem patet, $\delta$ iacere debere inter $-180^0$ et $+180^0$, atque $\epsilon$ inter 0 et $360^0$: sed sic quoque uterque angulus determinationem duplicem, adeoque tempus resultans quadruplicem admittere videtur. Attamen ex aequatione 5 art. 88 habemus
\[
\cos f\mathpunct{.}\BareRoot{rr'}=a(\cos g-\cos h)=2a\sin\tfrac12\delta\sin\tfrac12\epsilon:
\]
\noindent iam $\sin\tfrac12\epsilon$ necessario fit quantitas positiva, unde concludimus, $\cos f$ et $\sin\tfrac12\delta$ necessario eodem signo affectos esse, adeoque $\delta$ inter 0 et $180^0$ vel inter $-180^0$ et 0 accipiendum esse, prout $\cos f$ positivus fuerit vel negativus, i. e. prout motus heliocentricus $2f$ fuerit infra vel supra $180^0$. Ceterum sponte patet, pro $2f=180^0$ necessario esse debere $\delta=0$. Hoc itaque modo $\delta$ plene determinatus est. At determinatio anguli $\epsilon$ necessario ambigua manet, ita ut semper pro tempore \textit{duo} valores prodeant, quorum quis verus sit, nisi aliunde constet, decidi nequit. Ceterum ratio huius phaenomeni facile perspicitur: constat enim, per duo puncta data describi posse \textit{duas} ellipses diversas, quae ambae focum suum habeant in eodem puncto dato, simulque eundem semiaxem maiorem\footnote{Descripto e loco primo circulo radio $2a-r$, alioque radio $2a-r'$ e loco secundo, ellipseos focum alterum in intersectione horum circulorum iacere patet. Quare quum generaliter loquendo duae semper dentur intersectiones, duae ellipses diversae prodibunt.}; manifesto autem motus a loco primo ad secundum in his ellipsibus temporibus inaequalibus absolvetur.
